\documentclass[10pt]{beamer}
\usetheme{Madrid}
\usecolortheme{beaver}
\usepackage[utf8]{inputenc}
\usepackage[T1]{fontenc}
\usepackage[french]{babel}
\usepackage{hyperref}
\title{Soutenance du projet \textit{SecureChat}}
\subtitle{Réalisation d'un chat temps réel avec ASP.NET Core, EF Core, Identity et SignalR}
\author{Étudiant Ingénieur}
\date{\today}

\begin{document}

% Page de titre
\begin{frame}
  \titlepage
\end{frame}

% Plan
\begin{frame}{Plan (15 minutes)}
  \begin{enumerate}
    \item Contexte et objectifs
    \item Architecture globale et pile technologique
    \item Données et sécurité (EF Core + Identity)
    \item Messagerie temps réel (SignalR + Hub)
    \item Interfaces et API d'appoint
    \item Démonstration guidée et flux utilisateur
    \item Qualité, tests et pistes d'amélioration
  \end{enumerate}
\end{frame}

% Contexte et objectifs
\begin{frame}{Contexte et objectifs}
  \begin{itemize}
    \item Besoin: messagerie sécurisée pour utilisateurs authentifiés (DM + salons).
    \item Contraintes: simplicité de déploiement (SQLite), authentification intégrée, réactivité en temps réel.
    \item Cibles: preuve de concept robuste pour un cours d'architecture web moderne.
    \item Livrables principaux: application ASP.NET Core, base SQLite embarquée, hub SignalR, page unique de chat et API de recherche d'utilisateurs.
  \end{itemize}
\end{frame}

% Architecture globale
\begin{frame}{Architecture globale}
  \begin{block}{Couche serveur (Program.cs)}
    \begin{itemize}
      \item Enregistrement des services: EF Core (SQLite), Identity avec utilisateur personnalisé, Razor Pages, MVC, SignalR, \texttt{IUserIdProvider} custom.
      \item Pipeline: HTTPS, fichiers statiques, authentification/autorisation, routage Razor/MVC, hub chat exposé en \texttt{/hubs/chat}, endpoint \texttt{/whoami} pour debug.
    \end{itemize}
  \end{block}
  \begin{block}{Base de données}
    \begin{itemize}
      \item Connection string \texttt{DefaultConnection} pointant vers \texttt{app.db} (SQLite).
      \item Contexte unique \texttt{ApplicationDbContext} héritant d'Identity pour gérer tables + champs personnalisés.
    \end{itemize}
  \end{block}
\end{frame}

% Modèle de données
\begin{frame}{Données et sécurité (EF Core + Identity)}
  \begin{itemize}
    \item Modèle utilisateur \texttt{ApplicationUser} : \texttt{Email}, \texttt{PasswordHash} (Identity) + \texttt{DisplayName}, \texttt{AvatarUrl}.
    \item Contexte \texttt{ApplicationDbContext} : configuration des longueurs, migrations Identity prêtes via EF Core.
    \item Authentification: \texttt{[Authorize]} sur le hub SignalR et l'API de recherche \Rightarrow accès réservé aux comptes connectés.
    \item Avantage EF Core: un seul DbContext pour Identity et nos extensions, requêtes LINQ, migrations automatiques.
  \end{itemize}
\end{frame}

% SignalR
\begin{frame}{Messagerie temps réel (SignalR)}
  \begin{itemize}
    \item Hub \texttt{ChatHub} (dossier \texttt{Hubs/}) protégé par \texttt{[Authorize]}.
    \item Méthodes clés:
      \begin{itemize}
        \item \texttt{SendDirectByEmail}: résolution du destinataire par email + envoi ciblé \texttt{ReceiveDirect}.
        \item \texttt{JoinRoom}: ajout au groupe, diffusion d'un message "System" de bienvenue.
        \item \texttt{SendToRoom}: broadcast \texttt{ReceiveRoom} avec label émetteur et horodatage UTC.
      \end{itemize}
    \item \texttt{IUserIdProvider} personnalisé: mappe l'ID SignalR sur l'identité ASP.NET Core (NameIdentifier).
  \end{itemize}
\end{frame}

% UI
\begin{frame}{Interface utilisateur (Pages/Chat.cshtml)}
  \begin{itemize}
    \item Page Razor protégée par \texttt{[Authorize]}, trois blocs: messages directs, salons, journal temps réel.
    \item Client JS SignalR: connexion à \texttt{/hubs/chat}, handlers \texttt{ReceiveDirect}, \texttt{ReceiveRoom}, \texttt{System}.
    \item Autocomplete email: requêtes \texttt{/api/users/search?term=} avec debounce, remplit \texttt{datalist} pour choisir un destinataire.
    \item UX: nettoyage des champs après envoi, labels utilisateurs préférant \texttt{DisplayName} puis \texttt{Email}.
  \end{itemize}
\end{frame}

% API
\begin{frame}{API complémentaire (Controllers/UsersController.cs)}
  \begin{itemize}
    \item Route: \texttt{GET /api/users/search?term=...} sous \texttt{[Authorize]}.
    \item Filtre: préfixe insensible à la casse sur \texttt{Email} ou \texttt{DisplayName}, limité à 10 résultats.
    \item Payload: \texttt{Id}, \texttt{Email}, \texttt{DisplayName}. Sert à l'autocomplétion côté client.
  \end{itemize}
\end{frame}

% Démonstration
\begin{frame}{Démonstration guidée (5 min)}
  \begin{enumerate}
    \item Connexion/inscription via Identity (pages dans \texttt{Areas/Identity}).
    \item Accès \texttt{/Chat}: ouverture de la connexion SignalR.
    \item DM: saisie email, autocomplétion, envoi \Rightarrow événement \texttt{ReceiveDirect} sur le compte cible.
    \item Salon: rejoindre avec un identifiant, publication \Rightarrow \texttt{ReceiveRoom} pour tous les membres.
    \item Debug: endpoint \texttt{/whoami} pour vérifier l'ID utilisateur si besoin.
  \end{enumerate}
\end{frame}

% Qualité et tests
\begin{frame}{Qualité, performances et sécurité}
  \begin{itemize}
    \item Sécurité: tout le temps derrière \texttt{[Authorize]}, pas de données hors Identity; validation côté hub.
    \item Base légère: SQLite embarqué, idéal pour démo et tests rapides; peut être remplacé par SQL Server si nécessaire.
    \item Observabilité: logs ASP.NET Core + messages système dans les salons pour tracer les connexions.
    \item Robustesse SignalR: reconnexion automatique côté client, horodatage UTC pour cohérence.
  \end{itemize}
\end{frame}

% Perspectives
\begin{frame}{Pistes d'amélioration}
  \begin{itemize}
    \item Ajouter pièces jointes ou réactions dans les salons.
    \item Historiser les messages via une table EF Core dédiée.
    \item Gestion des statuts en ligne et présence (SignalR groups + cache distribué).
    \item Tests end-to-end (Playwright) et tests de hubs avec \texttt{Microsoft.AspNetCore.SignalR.Client} en mémoire.
  \end{itemize}
\end{frame}

% Conclusion
\begin{frame}{Conclusion}
  \begin{itemize}
    \item Projet illustre l'intégration cohérente EF Core + Identity + SignalR dans ASP.NET Core.
    \item Démo focalisée sur l'expérience utilisateur simple et sécurisée.
    \item Architecture modulaire: un DbContext, un hub, une page principale, une API d'appoint.
  \end{itemize}
  \vfill
  \begin{center}
    Merci pour votre attention \& questions
  \end{center}
\end{frame}

\end{document}
